\chapter{Kod źródłowy serwera WebSocket (server.py)}
\label{app:server}

Poniżej przedstawiono pełny kod pliku \texttt{broadcaster/server.py}. Serwer obsługuje tokeny, SSL, autoryzację ID, logowanie do pliku, powiadomienia Telegram oraz opcjonalny interfejs webowy (dashboard).

\lstset{language=Python, basicstyle=\ttfamily\footnotesize, breaklines=true, frame=single, numbers=left, numberstyle=\tiny, caption={}, captionpos=b}
\begin{lstlisting}
"""
Serwer WebSocket do przesyłania ramek CAN w czasie rzeczywistym.
Obsługuje opcjonalny token autoryzacyjny, autoryzację ID, logowanie do pliku,
powiadomienia Telegram oraz interfejs webowy (dashboard).
"""

import asyncio
import json
import logging
from datetime import datetime
from typing import Set, Optional
import ssl
import urllib.request
import urllib.parse
import os

import websockets
from websockets.server import WebSocketServerProtocol

\# aiohttp jest opcjonalne – jeśli nie ma, dashboard HTTP nie będzie dostępny
try:
    from aiohttp import web
    AIOHTTP\_AVAILABLE = True
except ImportError:
    AIOHTTP\_AVAILABLE = False
    web = None

logger = logging.getLogger(\_\_name\_\_)


class CANWebSocketServer:
    def \_\_init\_\_(self, host: str = "0.0.0.0", port: int = 8765,
                 token: Optional[str] = None, ssl\_context: Optional[ssl.SSLContext] = None):
        self.host = host
        self.port = port
        self.token = token
        self.ssl\_context = ssl\_context
        self.\_server = None
        self.\_clients: Set[WebSocketServerProtocol] = set()

        self.can\_interface = None
        self.allowed\_client\_ids = None
        self.incoming\_filter\_ids = None
        self.log\_to\_file = False
        self.log\_file\_path = "remote\_operations.log"

        self.telegram\_token = None
        self.telegram\_chat\_id = None

        self.enable\_http = False
        self.http\_port = 8080
        self.http\_app = None
        self.http\_runner = None

        self.\_running = False

    def \_send\_telegram(self, message):
        if not self.telegram\_token or not self.telegram\_chat\_id:
            return
        try:
            url = f"https://api.telegram.org/bot\{self.telegram\_token\}/sendMessage"
            data = urllib.parse.urlencode(\{
                "chat\_id": self.telegram\_chat\_id,
                "text": f"🚗 CAN Simulator\textbackslash n\{message\}",
                "parse\_mode": "HTML"
            \}).encode("utf-8")
            req = urllib.request.Request(url, data=data)
            with urllib.request.urlopen(req, timeout=5) as resp:
                if resp.getcode() != 200:
                    logger.warning(f"Telegram odpowiedział kodem \{resp.getcode()\}")
        except Exception as e:
            logger.error(f"Błąd wysyłania do Telegram: \{e\}")

    async def \_handler(self, websocket: WebSocketServerProtocol):
        logger.info(f"Nowy klient: \{websocket.remote\_address\}")

        if self.token:
            try:
                msg = await asyncio.wait\_for(websocket.recv(), timeout=5.0)
                if msg != self.token:
                    logger.warning(f"Nieprawidłowy token od \{websocket.remote\_address\}")
                    await websocket.close(1008, "Invalid token")
                    return
                logger.info(f"Token zaakceptowany od \{websocket.remote\_address\}")
            except asyncio.TimeoutError:
                logger.warning(f"Timeout tokena od \{websocket.remote\_address\}")
                await websocket.close(1008, "Token timeout")
                return

        if self.allowed\_client\_ids is not None:
            try:
                msg = await asyncio.wait\_for(websocket.recv(), timeout=5.0)
                data = json.loads(msg)
                client\_ids = set(data.get("allowed\_ids", []))
                if not client\_ids.issubset(self.allowed\_client\_ids):
                    logger.warning(f"Klient \{websocket.remote\_address\} próbował użyć niedozwolonych ID")
                    await websocket.close(1008, "Forbidden IDs")
                    return
                websocket.client\_allowed\_ids = client\_ids
                logger.info(f"Klient \{websocket.remote\_address\} autoryzowany z ID: \{client\_ids\}")
            except (asyncio.TimeoutError, json.JSONDecodeError):
                logger.warning(f"Błąd autoryzacji ID od \{websocket.remote\_address\}")
                await websocket.close(1008, "Invalid ID list")
                return

        self.\_clients.add(websocket)
        logger.info(f"Klient dodany, łącznie: \{len(self.\_clients)\}")
        self.\_send\_telegram(f"🟢 Klient połączony: \{websocket.remote\_address\}")

        try:
            async for message in websocket:
                await self.\_handle\_incoming\_frame(websocket, message)
        except websockets.exceptions.ConnectionClosed:
            logger.info(f"Połączenie zamknięte przez klienta \{websocket.remote\_address\}")
        finally:
            self.\_clients.remove(websocket)
            logger.info(f"Klient rozłączony, pozostało: \{len(self.\_clients)\}")
            self.\_send\_telegram(f"🔴 Klient rozłączony: \{websocket.remote\_address\}")

    async def \_handle\_incoming\_frame(self, websocket, message):
        if not self.incoming\_filter\_ids:
            logger.debug("Brak filtru przychodzącego – ramka odrzucona.")
            return
        try:
            data = json.loads(message)
            can\_id = int(data['id'], 16) if isinstance(data['id'], str) else data['id']
            if can\_id not in self.incoming\_filter\_ids:
                logger.debug(f"ID 0x\{can\_id:X\} nie przechodzi filtra przychodzącego.")
                return
            payload = bytes(data['data'])
            is\_extended = data.get('is\_extended', False)
            if self.can\_interface and self.can\_interface.connected:
                success, msg = self.can\_interface.send\_frame(can\_id, payload, is\_extended)
                if success:
                    logger.info(f"Wysłano na CAN: ID=0x\{can\_id:X\}, data=\{payload.hex().upper()\}")
                    if self.log\_to\_file:
                        self.\_log\_to\_file(f"RX from \{websocket.remote\_address\}: \{message\}")
                else:
                    logger.error(f"Błąd wysyłania na CAN: \{msg\}")
                    self.\_send\_telegram(f"❌ Błąd wysyłania na CAN: \{msg\}")
            else:
                logger.warning("CAN niepodłączony – ramka odrzucona.")
        except Exception as e:
            logger.error(f"Błąd przetwarzania ramki od klienta: \{e\}")

    def \_create\_http\_app(self):
        if not AIOHTTP\_AVAILABLE:
            return None
        app = web.Application()
        web\_dir = os.path.join(os.path.dirname(os.path.dirname(\_\_file\_\_)), 'web\_dashboard')
        if os.path.isdir(web\_dir):
            app.router.add\_static('/', web\_dir, show\_index=True)
        return app

    async def start(self):
        if self.\_running:
            return
        logger.info(f"Uruchamianie serwera WebSocket na \{self.host\}:\{self.port\} (SSL: \{self.ssl\_context is not None\})")
        self.\_server = await websockets.serve(
            self.\_handler, self.host, self.port, ssl=self.ssl\_context
        )
        self.\_running = True

        if self.enable\_http:
            self.http\_app = self.\_create\_http\_app()
            if self.http\_app:
                runner = web.AppRunner(self.http\_app)
                await runner.setup()
                site = web.TCPSite(runner, self.host, self.http\_port)
                await site.start()
                self.http\_runner = runner
                logger.info(f"Interfejs webowy dostępny na http://\{self.host\}:\{self.http\_port\}")

    async def stop(self):
        if not self.\_running:
            return
        logger.info("Zatrzymywanie serwera...")
        self.\_server.close()
        await self.\_server.wait\_closed()
        for client in list(self.\_clients):
            await client.close()
        self.\_clients.clear()

        if self.http\_runner:
            await self.http\_runner.cleanup()
            self.http\_runner = None

        self.\_running = False
        logger.info("Serwer zatrzymany")

    def \_log\_to\_file(self, msg):
        try:
            with open(self.log\_file\_path, 'a', encoding='utf-8') as f:
                f.write(f"\{datetime.now().isoformat()\} \{msg\}\textbackslash n")
        except Exception as e:
            logger.error(f"Błąd zapisu do pliku logu: \{e\}")

    def broadcast\_frame(self, frame: dict):
        if not self.\_running or not self.\_clients:
            return
        message = json.dumps(frame)
        for client in list(self.\_clients):
            try:
                asyncio.create\_task(client.send(message))
            except Exception as e:
                logger.error(f"Błąd wysyłania: \{e\}")

    async def broadcast\_frame\_async(self, frame: dict):
        if not self.\_running or not self.\_clients:
            return
        message = json.dumps(frame)

        if self.log\_to\_file:
            self.\_log\_to\_file(f"TX to clients: \{message\}")

        for client in list(self.\_clients):
            try:
                await client.send(message)
            except Exception as e:
                logger.error(f"Błąd wysyłania async: \{e\}")

    @property
    def client\_count(self) -> int:
        return len(self.\_clients)

    @property
    def is\_running(self) -> bool:
        return self.\_running
\end{lstlisting}
